\documentclass[12pt, babellanguage=british, titlespace, defaulttheorems]{altarticle}

\title{Documentation for the \texttt{altarticle} class}
\date{}

\newcommand{\optionline}[2][]{\par\medskip\hspace*{-.75cm}\texttt{#2}\hfill\if\relax\detokenize{#1}\relax\else\texttt{default:~#1}\fi\par}
\newcommand{\aargt}[1]{\ensuremath{\langle}\textrm{{\emph{#1}}}\ensuremath{\rangle}}
\newcommand{\argt}[1]{\texttt{\{}\ensuremath{\langle}\textrm{{\emph{#1}}}\ensuremath{\rangle}\texttt{\}}}
\newcommand{\optargt}[1]{\texttt{[}\ensuremath{\langle}\textrm{\emph{#1}}\ensuremath{\rangle}\texttt{]}}
\newcommand{\commandline}[1][]{\def\argI{#1}\ccommandline}
\newcommand{\ccommandline}[2][]{\par\medskip\hspace*{-.75cm}\texttt{\textbackslash#2}\argI{}\hfill\if\relax\detokenize{#1}\relax\else\texttt{default:~#1}\fi\par}
\newcommand{\pcommandline}[1][]{\def\argI{#1}\pccommandline}
\newcommand{\pccommandline}[2][]{\par\hspace*{-.75cm}\hphantom{\texttt{\textbackslash#2}}\argI{}\hfill\if\relax\detokenize{#1}\relax\else\texttt{default:~#1}\fi\par}
\newcommand{\envline}[1][]{\def\argI{#1}\eenvline}
\newcommand{\eenvline}[2][]{\par\medskip\hspace*{-.75cm}\texttt{\textbackslash begin\{#2\}}\argI{}~...~\texttt{\textbackslash end\{#2\}}\hfill\if\relax\detokenize{#1}\relax\else\texttt{default:~#1}\fi\par}

\begin{document}
\maketitle
\tableofcontents
\bigskip
This class is based on the \texttt{article} class, and the commands common to the \texttt{alt} classes are available in the documentation for \texttt{altcommands}. It is made for documents either in French or in English.
\section{Options}
On top of the options of the article class, some options are available to customise the apprearance of the documents.

\optionline[margin=2.5cm]{geometryargs}This option controls the argument passed to the package \texttt{geometry}. If multiple arguments are passed, they put between brackets.
\optionline[french]{babellanguage}This options defines the languages processed by \texttt{babel}. When the language \texttt{french} is not loaded, all of the commands are typeset in English. They are otherwise typset in French. To have the commands loaded in English and load the French language in the document, \texttt{\textbackslash PassOptionsToPackage\{french\}\{babel\}} can be used before or after loading the class.
\optionline[none]{bibsorting}This option defines the sorting method used by \texttt{biblatex}. It has no effect if the option \texttt{bibliography} is not used.
\optionline[false]{unslantgreek}This option allows to have unslanted Greek letters in maths. When the \texttt{unslantgreek} command is not used, uppercase Greek letters are in italics and upright versions can be obtained by adding \texttt{up} at the beginning of the command, like \texttt{\textbackslash upDelta} or \texttt{\textbackslash upGamma}.
\optionline[false]{unslantlatin}This option allows to have unslanted lowercase Latin letters in maths.
\optionline[false]{unslantLatin}This option allows to have unslanted uppercase Latin letters in maths.
\optionline[false]{defaulttheorems}This options defines some theorems that can be used throughout the document. The ones created by this options are defined by
if the French language is loaded for \texttt{babel}
\begin{verbatim}
\createtheorem{axiome}{Axiome}{axi}
              {axiome}{section}{true}
\createtheorem{corollaire}{Corollaire}{cor}
              {corollaire}{section}{true}
\createtheorem{definition}{Définition}{def}
              {définition}{section}{true}
\createtheorem{lemme}{Lemme}{lm}
              {lemme}{section}{true}
\createtheorem{notation}{Notation}{not}
              {notation}{section}{true}
\createtheorem{proposition}{Proposition}{prop}
              {proposition}{section}{true}
\createtheorem{remarque}{Remarque}{rem}
              {remarque}{section}{true}
\createtheorem{theoreme}{Théorème}{thm}
              {théorème}{section}{true}
\end{verbatim}
and
\begin{verbatim}
\createtheorem{axiom}{Axiom}{axi}
              {axiom}{section}{true}
\createtheorem{corollary}{Corollary}{cor}
              {corollary}{section}{true}
\createtheorem{definition}{Definition}{def}
              {definition}{section}{true}
\createtheorem{lemma}{Lemma}{lm}
              {lemma}{section}{true}
\createtheorem{notation}{Notation}{not}
              {notation}{section}{true}
\createtheorem{proposition}{Proposition}{prop}
              {proposition}{section}{true}
\createtheorem{remark}{Remark}{rem}
              {remark}{section}{true}
\createtheorem{theorem}{Theorem}{thm}
              {theorem}{section}{true}
\end{verbatim}
otherwise
\optionline[false]{bibliography}When this option is used, the package \texttt{biblatex} is loaded and modifications are made to its default behaviour. It also creates two files for bibliography printing.
\optionline[false]{index}This option provides commands to create an index out of the named theorems created by this class. It creates a \texttt{Python} script that will be called upon creating the index. To create the index, the \texttt{-shell-escape} compile parameter must be used, and \texttt{Python} must be acceddible via the command \texttt{python3}.
\optionline[false]{titlespace}This command allows to add spacing around the \texttt{\textbackslash maketitle} command.
\optionline[false]{draft}This option overrides, and passes the option to the \texttt{article} class. In order to use the draft option of the \texttt{article} class without the one defined by this package, the command \texttt{\textbackslash PassOptionsToClass\{draft\}\{article\}} must be used before loading the class.
\section{Loaded packages and class}
Here are the loaded packages and classes, as well as the default parameter and redefinitions made:
\begin{verbatim}
\LoadClass[a4paper]{article}
\usepackage{altcommands}
\usepackage{geometry}
\usepackage{framed}
\usepackage{enumitem}
\usepackage{multicol}
\usepackage{tabularray}

% Remove space around multicols
\BeforeBeginEnvironment{multicols}
    {\par\vphantom{Ip}\vspace{-\baselineskip}}
\AfterEndEnvironment{multicols}
    {\par\vphantom{Ip}\vspace{-\baselineskip}}

% tabularray options
\SetTblrInner{rowsep=1pt}
\DefTblrTemplate{caption-tag}{default}{}
\DefTblrTemplate{caption-sep}{default}{}
\DefTblrTemplate{caption-text}{default}{}
\DefTblrTemplate{contfoot-text}{default}{}
\DefTblrTemplate{conthead-text}{default}{}
\end{verbatim}

\null

Lists are also (re)defined by the class.
\commandline[]{itembn\argt{number}}This commands prints the number with the \texttt{oldstylenums} style between parentheses, like \itemnb1 and \itemnb{125}. It is used for ordered lists enumeration.
\envline{propenumerate}This is a variant of the \texttt{enumerate} environment with spacing adapted to the proof environment.
\envline{propitemize}This is a variant of the \texttt{itemize} environment with spacing adapted to the proof environment.
\envline{prvenumerate}This is a variant of the \texttt{enumerate} environment with no spacing.
\envline{prvitemize}This is a variant of the \texttt{itemize} environment with no spacing.

\null

\commandline{cleardoublepage}This command is redefined to behave as if the document was twosided. The original behaviour of the command can be obtained by using \texttt{\textbackslash oldcleardoublepage}.
\section{Draft mode}
In draft mode, a header is created indicating the time at which the document was compiled. It is done by defining the \texttt{pagestyle} to \texttt{versionheader}. The colour used by the header is the colour \texttt{headercolour}, which can be redefined. This header can also be removed by changing the \texttt{pagestyle}. To remove the version information on the page containing the title, it is necessary to use \texttt{\textbackslash apptocmd\{\textbackslash maketitle\}\{\textbackslash thispagestyle\{...\}\}\{\}\{\}}.
\section{Theorems}
\subsection{Theorem commands}
This class implements its own version of theorems. The commands linked to theorems are described below.
\commandline[\argt{command}\argt{name}\argt{labelname}]{createtheorem}
\pcommandline[\argt{refname}\argt{parentcounter}\argt{sharedcounter}]{createtheorem}This command creates a theorem that can be used with the command given as argument.

The \aargt{name} argument corresponds to the text that will be displayed upon calling the theorem.

The \aargt{labelname} corresponds to the name given as prefix for references, when calling \texttt{\textbackslash autoref} with the theorem.

The \aargt{parentcounter} argument defines the parent counter to be used as the counter of the theorem. It can be left empty for the counter to be global within the document.

The \aargt{sharedcounter} bolean option, that behaves as false if \texttt{true} is not given as input, indicates if the counter of the theorem is to be shared with the other theorems having the same parent counter.

For instance,
\begin{verbatim}
\createtheorem{definition}{Définition}{def}
              {définition}{section}{true}
\end{verbatim}
creates a theorem that can be used with \texttt{\textbackslash definition}, and with a counter that resets at every section, and with its counter incremented with the other theorems with the same shared counter within the sections.
\commandline[\optargt{ref}\optargt{complement}\optargt{altname}\optargt{hidename}\argt{name}][{[][][][]}]{\emph{theorem}}
This command, with \texttt{\aargt{theorem}} being replaced by the command name provided in the previous \texttt{\textbackslash createtheorem} command, creates a numbered theorem with name \aargt{name}. If the name contains commands, or non-ascii characters, an \aargt{ref} might have to be given to prevent reference errors or warnings. When starting the \aargt{name} with a character outside of range \texttt{a-zA-Z}, it might be needed to put it within brackets if an index is created.

A label to this theorem is created whenever \aargt{ref} or \aargt{name} is not empty, with \aargt{ref} prefered over \aargt{name}. The reference label name starts with \aargt{labelname} provided to the \texttt{\textbackslash createtheorem} command.

The \aargt{complement} argument corresponds to text that will be added before, if French is loaded, or after otherwise whenever a reference to the theorem is made.

The \aargt{altname} argument corresponds to a variant of the \aargt{name} that is to be used when math is not in bold mode. If none is provided, the \aargt{name} is used instead.

The \aargt{hidename} argument indicates whether the \aargt{name} of the theorem has to be displayed in the references or if only the number is to be displayed.
\commandline[*\optargt{ref}\optargt{complement}\optargt{altname}\optargt{hidename}\argt{name}][{[][][][]}]{\emph{theorem}}
This command, is the same as the unstarred one, except that it creates an unnumbered theorem.

The \aargt{hidename} is here for compatibility reasons with the non-starred theorem command, but should not be used with starred theorems.
\commandline[\argt{thmlabel}]{thmref}This implements a reference system that is to be used upon making a reference to a theorem. If it is used for other elements than theorems, it will behave like \texttt{\textbackslash autoref}.
\commandline[\argt{thmlabel}]{thmnameref}This command is similar to \texttt{\textbackslash thmref}, except that it forces the display of the name of the theorem, even if \aargt{hidename} is set to \texttt{true}.
\commandline[][\textbackslash thmnametrue]{thmname}This corresponds to a boolean that indicates whether the name of the theorem is to be displayed when creating them. It can be changed with \texttt{\textbackslash thmnamefalse} and \texttt{\textbackslash thmnametrue}.
\envline[\optargt{extra}\argt{thmref}][{[]}]{proof}This environment allows to create proofs, with \aargt{thmref} being the reference to a theorem created in the document. The text \aargt{extra} is to be added right after the reference to the theorem and is initially empty.
\envline[\optargt{extra}\argt{thmref}][{[]}]{preuve}This environment is the same as the \texttt{proof} one.
\subsection{Examples}
The code
\begin{verbatim}
\theorem{Euclid}For every prime $p$, there is a prime $p'>p$.
    In particular, there are infinitely many primes.
\begin{proof}{thm Euclid}
    Lorem ipsum \dots
\end{proof}
\theorem[euclid theorem]['s theorem]{Euclid}For every prime $p$, there
    is a prime $p'>p$. In particular, there are infinitely many primes.
\begin{proof}{thm euclid theorem}
    Lorem ipsum \dots
\end{proof}
\theorem[euclid 2][]{Euclid}For every prime $p$, there is a prime
    $p'>p$. In particular, there are infinitely many primes.
\begin{proof}{thm euclid 2}
    Lorem ipsum \dots
\end{proof}
\theorem[euclid 3][][normal Euclid]{bold Euclid}For every prime $p$,
    there is a prime $p'>p$. In particular, there are infinitely many
    primes.
\begin{proof}{thm euclid 3}
    Lorem ipsum \dots
\end{proof}
With \thmref{thm euclid 3} and {\boldmath\thmref{thm euclid 3}}.
\theorem[euclid 4][][][true]{not Euclid}For every prime $p$, there is a
    prime $p'>p$. In particular, there are infinitely many primes.
\begin{proof}{thm euclid 4}
    Lorem ipsum \dots
\end{proof}
\theorem*[euclid star theorem]['s theorem]{Euclid}For every prime $p$,
    there is a prime $p'>p$. In particular, there are infinitely many
    primes.
\begin{proof}[ -- starred version]{thm euclid star theorem}
    Lorem ipsum \dots
\end{proof}
\end{verbatim}
produces
\theorem{Euclid}For every prime $p$, there is a prime $p'>p$.In particular, there are infinitely many primes.
\begin{proof}{thm Euclid}
    Lorem ipsum \dots
\end{proof}
\theorem[euclid theorem]['s theorem]{Euclid}For every prime $p$, there is a prime $p'>p$. In particular, there are infinitely many primes.
\begin{proof}{thm euclid theorem}
    Lorem ipsum \dots
\end{proof}
\theorem[euclid 2][]{Euclid}For every prime $p$, there is a prime $p'>p$. In particular, there are infinitely many primes.
\begin{proof}{thm euclid 2}
    Lorem ipsum \dots
\end{proof}
\theorem[euclid 3][][normal Euclid]{bold Euclid}For every prime $p$, there is a prime $p'>p$. In particular, there are infinitely many primes.
\begin{proof}{thm euclid 3}
    Lorem ipsum \dots
\end{proof}
With \thmref{thm euclid 3} and {\boldmath\thmref{thm euclid 3}}.
\theorem[euclid 4][][][true]{not Euclid}For every prime $p$, there is a prime $p'>p$. In particular, there are infinitely many primes.
\begin{proof}{thm euclid 4}
    Lorem ipsum \dots
\end{proof}
\theorem*[euclid star theorem]['s theorem]{Euclid}For every prime $p$, there is a prime $p'>p$. In particular, there are infinitely many primes.
\begin{proof}[ -- starred version]{thm euclid star theorem}
    Lorem ipsum \dots
\end{proof}
\subsection{Index}
An index of all of the theorems and definitions with names in the document can be created. To create this index, \texttt{Python3} must be installed and the bash command \texttt{python3} should call the python executer. The \texttt{-shell-escape} compile option must be used to create the index.
\commandline{theoremtable}This command creates the index of the theorems and definitions. To modified the theorems inside of the index, it is possible to modify the command, which is implemented as
\begin{verbatim}
\newcommand\theoremtable{%
    \input{|python3 sorttables.py "\jobname"}%
    \section*{Index%
        \@mkboth{%
           \noexpand\expandafter\noexpand\MakeUppercase Index}
           {\noexpand\expandafter\noexpand\MakeUppercase Index}}%
    \addcontentsline{toc}{section}{\protect\numberline{}Index}%
    \begin{multicols}{2}\raggedcolumns%
    \@starttoc{totheoreme}%
    \@starttoc{totheorem}%
    \@starttoc{todefinition}%
    \end{multicols}%
    \def\toc{false}%
}
\end{verbatim}
It is also possible to manually add elements to the index, as for the table of contents, by adding the line \texttt{\textbackslash addcontentsline\{todefinition\}\{thm\}}\argt{contentname} after creating a theorem where \aargt{contentname} is the name to be displayed in the index.
\section{Captions}
Captions are redefined by the class to have automatic labels.
\commandline[\optargt{refname}\argt{name}][{[]}]{caption}This command creates text with the reference and also creates a reference with name <\aargt{float} \aargt{name}> where \aargt{float} is the name of the environment display in the caption like figure, table (tableau in French), \dots

The \aargt{refname} argument is used, if provided, to create the name of the reference if for instance the \aargt{name} contains some characters that are not comparible with references. The reference is then <\aargt{float} \aargt{refname}>
\commandline[\optargt{refname}\argt{name}][{[]}]{subcaption}This command is the same as \texttt{\textbackslash caption}, except that it creates a caption for a sub-object. The reference is to be made with the name <sub\aargt{float} \aargt{name}> or <sub\aargt{float} \aargt{refname}>.
\section{Default lengths}
The \texttt{altarticle} class redefines some lengths with
\begin{verbatim}
\AtBeginDocument{%
    \setlength{\belowdisplayskip}{0pt}%
    \setlength{\belowdisplayshortskip}{0pt}%
    \setlength{\abovedisplayskip}{0pt}%
    \setlength{\abovedisplayshortskip}{0pt}%
    \setlength{\parindent}{0pt}%
    \setlength{\topsep}{0pt}%
    \setlength{\partopsep}{0pt}%
    \setlength\multicolsep{0pt}%
}
\setlength{\dotsep}{3.5pt}
\setlength{\dotlift}{0.5pt}
\setlength{\matsep}{3.5pt}
\setlength{\matmin}{15pt}
\end{verbatim}
\end{document}